% problem_id: S001-lemma-cd-1-canonical
% source_id: S001 (Stacks Project)
% source_file: algebraization.tex
% statement_locator: algebraization.tex lines 7128--7141
% proof_locator: algebraization.tex lines 7143--7260
% env: lemma
% id_note: label 在本来源内唯一
% pairing: tight (verified)
% status: complete (statement + author proof verbatim + reason + locators)
%
\begin{lemma}
\label{lemma-cd-1-canonical}
In Situation \ref{situation-algebraize} let $(\mathcal{F}_n)$
be an object of $\textit{Coh}(U, I\mathcal{O}_U)$. Assume
\begin{enumerate}
\item $A$ has a dualizing complex and $\text{cd}(A, I) = 1$,
\item $(\mathcal{F}_n)$ is pro-isomorphic to an inverse system
$(\mathcal{F}_n'')$ of coherent $\mathcal{O}_U$-modules such that
$\text{depth}(\mathcal{F}''_{n, y}) + \delta^Y_Z(y) \geq 3$
for all $y \in U \cap Y$.
\end{enumerate}
Then $(\mathcal{F}_n)$ extends canonically to $X$, see
Definition \ref{definition-canonically-algebraizable}.
\end{lemma}

\begin{proof}
We will check hypotheses (a), (b), and (c) of Lemma \ref{lemma-when-done}.
Before we start, let us point out that the modules
$H^0(U, \mathcal{F}''_n)$ and $H^1(U, \mathcal{F}''_n)$
are finite $A$-modules for all $n$ by
Local Cohomology, Lemma \ref{local-cohomology-lemma-finiteness-Rjstar}.

\medskip\noindent
Observe that for each $p \geq 0$
the limit topology on $\lim H^p(U, \mathcal{F}_n)$
is the $I$-adic topology by Lemma \ref{lemma-topology-I-adic}.
In particular, hypothesis (b) holds.

\medskip\noindent
We know that $M = \lim H^0(U, \mathcal{F}_n)$ is an $A$-module whose
limit topology is the $I$-adic topology. Thus, given $n$, the module
$M/I^nM$ is a subquotient of $H^0(U, \mathcal{F}_N)$ for some $N \gg n$.
Since the inverse system $\{H^0(U, \mathcal{F}_N)\}$ is pro-isomorphic to an
inverse system of finite $A$-modules, namely $\{H^0(U, \mathcal{F}''_N)\}$,
we conclude that $M/I^nM$ is finite. It follows that $M$ is finite, see
Algebra, Lemma \ref{algebra-lemma-finite-over-complete-ring}.
In particular hypothesis (c) holds.

\medskip\noindent
For each $n \geq 0$ let us write $Ob_n = \lim_N H^1(U, I^n\mathcal{F}_N)$.
A special case is $Ob = Ob_0 = \lim_N H^1(U, \mathcal{F}_N)$.
Arguing exactly as in the previous paragraph we find that $Ob$
is a finite $A$-module. (In fact, we also know that $Ob/I Ob$ is annihilated
by a power of $\mathfrak a$, but it seems somewhat difficult to use this.)

\medskip\noindent
We set $\mathcal{F} = \lim \mathcal{F}_n$, we pick generators
$f_1, \ldots, f_r$ of $I$, we pick $c \geq 1$, and we choose
$\Phi_\mathcal{F}$ as in Lemma \ref{lemma-cd-is-one-for-system}.
We will use the results of Lemma \ref{lemma-properties-system}
without further mention. In particular, for each $n \geq 1$ there are maps
$$
\delta_n :
H^0(U, \mathcal{F}_n)
\longrightarrow
H^1(U, I^n\mathcal{F})
\longrightarrow
Ob_n
$$
The first comes from the short exact sequence
$0 \to I^n\mathcal{F} \to \mathcal{F} \to \mathcal{F}_n \to 0$
and the second from $I^n\mathcal{F} = \lim I^n\mathcal{F}_N$.
We will later use that if $\delta_n(s) = 0$ for $s \in H^0(U, \mathcal{F}_n)$
then we can for each $n' \geq n$ find $s' \in H^0(U, \mathcal{F}_{n'})$
mapping to $s$.
Observe that there are commutative diagrams
$$
\xymatrix{
H^0(U, \mathcal{F}_{nc}) \ar[r] \ar[dd] &
H^1(U, I^{nc}\mathcal{F}) \ar[dd] \ar[rd]^{\Phi_\mathcal{F}} \\
& &
\bigoplus_{e_1 + \ldots + e_r = n}
H^1(U, \mathcal{F}) \cdot T_1^{e_1} \ldots T_r^{e_r} \ar[ld] \\
H^0(U, \mathcal{F}_n) \ar[r] &
H^1(U, I^n\mathcal{F})
}
$$
We conclude that the obstruction map
$H^0(U, \mathcal{F}_n) \to Ob_n$
sends the image of
$H^0(U, \mathcal{F}_{nc}) \to H^0(U, \mathcal{F}_n)$
into the submodule
$$
Ob'_n =
\Im\left(
\bigoplus\nolimits_{e_1 + \ldots + e_r = n}
Ob \cdot T_1^{e_1} \ldots T_r^{e_r} \to Ob_n
\right)
$$
where on the summand $Ob \cdot T_1^{e_1} \ldots T_r^{e_r}$
we use the map on cohomology coming from the reductions modulo
powers of $I$ of the multiplication map
$f_1^{e_1} \ldots f_r^{e_r} : \mathcal{F} \to I^n\mathcal{F}$.
By construction
$$
\bigoplus\nolimits_{n \geq 0} Ob'_n
$$
is a finite graded module over the Rees algebra $\bigoplus_{n \geq 0} I^n$.
For each $n$ we set
$$
M_n = \{s \in H^0(U, \mathcal{F}_n) \mid \delta_n(s) \in Ob'_n\}
$$
Observe that $\{M_n\}$ is an inverse system and that
$f_j : \mathcal{F}_n \to \mathcal{F}_{n + 1}$ on global
sections maps $M_n$ into $M_{n + 1}$.
By exactly the same argument as in the proof of
Cohomology, Lemma \ref{cohomology-lemma-ML-general}
we find that $\{M_n\}$ is ML. Namely, because the Rees algebra
is Noetherian we can choose a finite number of homogeneous generators
of the form $\delta_{n_j}(z_j)$ with $z_j \in M_{n_j}$ for the graded submodule
$\bigoplus_{n \geq 0} \Im(M_n \to Ob'_n)$.
Then if $k = \max(n_j)$ we find that for $n \geq k$
and any $z \in M_n$ we can find $a_j \in I^{n - n_j}$ such that
$z - \sum a_j z_j$ is in the kernel of $\delta_n$
and hence in the image of $M_{n'}$ for all $n' \geq n$
(because the vanishing of $\delta_n$ means that we can
lift $z - \sum a_j z_j$ to an element $z' \in H^0(U, \mathcal{F}_{n'c})$
for all $n' \ge n$ and then the image of $z'$ in $H^0(U, \mathcal{F}_{n'})$
is in $M_{n'}$ by what we proved above).
Thus $\Im(M_n \to M_{n - k}) = \Im(M_{n'} \to M_{n - k})$
for all $n' \geq n$.

\medskip\noindent
Choose $n$. By the Mittag-Leffler property of $\{M_n\}$ we just established
we can find an $n' \geq n$ such that the image of $M_{n'} \to M_n$
is the same as the image of $M' \to M_n$. By the above we see that
the image of $M' \to M_n$ contains the image of
$H^0(U, \mathcal{F}_{n'c}) \to H^0(U, \mathcal{F}_n)$.
Thus we see that $\{M_n\}$ and $\{H^0(U, \mathcal{F}_n)\}$
are pro-isomorphic. Therefore $\{H^0(U, \mathcal{F}_n)\}$
has ML and we finally conclude that hypothesis (a) holds.
This concludes the proof.
\end{proof}
